\section{The Night of the Denials}

\subsection{Supper, sword, and prophecy}

The last night begins, in John, with a protest. When Jesus kneels to wash
the disciples' feet, Peter resists---``Lord, dost thou wash my feet?\ldots\
Thou shalt never wash my feet''---and, told ``If I wash thee not, thou shalt
have no part with me,'' swings to the other extreme: ``Lord, not only my
feet, but also my hands and my head'' (Jn 13:6--9). The same evening all
four Gospels record the boast and the prophecy. In Mark: ``Although all
shall be scandalized in thee, yet not I,'' answered by ``today, even in this
night, before the cock crow twice, thou shalt deny me thrice,'' and the
more vehement reply, ``Although I should die together with thee, I will not
deny thee'' (Mk 14:29--31). In the garden it is John who names the swordsman:
``Then Simon Peter, having a sword, drew it, and struck the servant of the
high priest, and cut off his right ear. And the name of the servant was
Malchus,'' drawing the command, ``Put up thy sword into the scabbard. The
chalice which my Father hath given me, shall I not drink it?''
(Jn 18:10--11).

\subsection{Luke 22:31--32: the second primacy text}

Luke frames the prediction with two verses that carry more doctrinal weight
than any other saying about Peter outside Matthew 16, and that are usually
quoted without the feature that makes them interesting:

\begin{studyframe}\small
``And the Lord said: Simon, Simon, behold Satan hath desired to have you,
that he may sift you as wheat: But I have prayed for thee, that thy faith
fail not: and thou, being once converted, confirm thy brethren.''
(Lk 22:31--32)
\end{studyframe}

In English ``you'' hides the pivot. The Greek is explicit: Satan has demanded
\textit{hymas}, all of you, plural---the whole group is the target---but ``I
have prayed \textit{peri sou},'' concerning thee, singular, ``that
\textit{he pistis sou}, thy faith, may not fail''; and the charge that
follows, \textit{sterison tous adelphous sou}, ``confirm thy brethren,'' is
singular imperative to a man with brothers to strengthen. The sentence moves
from the many to the one and back to the many, and the one is named twice at
the start. Whatever it means, it does not mean that Peter is simply one of
the sifted.

Three things about the saying constrain any use of it.

\textbf{It is a prayer, not a guarantee of behaviour.} Luke places it
immediately before the boast and the prediction of the denial. Within thirty
verses Peter denies Jesus three times. The evangelist evidently did not think
the prayer meant that Peter would not fall; whatever ``that thy faith fail
not'' secures, it is compatible in Luke's own narrative with a catastrophic
public failure of nerve days later.

\textbf{It presupposes the fall.} ``Thou, being once converted''---the Greek
verb is \textit{epistrepho}, to turn back---builds the recovery into the
commission. The task of strengthening the others is given to a man on the
express assumption that he will first have to be turned around himself. Read
as an ecclesiological text it is remarkable for that reason: the confirming
of the brethren is entrusted precisely to the one who is about to fail
publicly and be restored.

\textbf{It says nothing about anyone after Peter.} Like Matthew 16, this is a
saying to a named man in a named crisis. The extension to his successors is a
doctrinal development, and the First Vatican Council makes exactly that
extension explicit, quoting the verse in its fourth chapter as ``the divine
promise made by our Lord and Saviour to the prince of his disciples'' and
inferring from it that the see of Peter ``remains ever untainted by any
error.'' A reader is entitled to think that inference sound or unsound; a
reader is not entitled to think that the verse states it. The council's own
next sentence, in the same chapter, is a limitation that is quoted far less
often: the Holy Spirit was promised to Peter's successors not that they might
``disclose new doctrine'' by his revelation, but that with his assistance
they might ``devoutly guard and faithfully expound'' the revelation handed
down through the apostles.

\subsection{The denials}

The denials themselves are among the best-attested scenes in the Gospels:
all four evangelists place Peter in the high priest's courtyard that night,
warming himself at the fire (``a fire of coals,'' Jn 18:18), and all four
record three denials before cockcrow. The witnesses vary in the
details---Mark alone counts two cockcrows; the challengers shift among
maidservants, bystanders, and, in John, a kinsman of Malchus; only Luke has
the terrible detail of the glance: ``And immediately, as he was yet speaking,
the cock crew. And the Lord turning looked on Peter'' (Lk 22:60--61).

The variations are the ordinary signature of independent tellings of one
remembered event. The event itself is as secure as narrative evidence
allows, and the reason deserves to be stated carefully rather than waved at.
It is not simply that the story is embarrassing; a religious movement can
tolerate a great deal of embarrassment about its founders' early failures if
the failure is safely past. It is that this particular failure is
load-bearing in the opposite direction. The four Gospels were written in
communities that treated Peter as the first of the apostles; three of them
report a saying or scene that gives him a distinctive commission; and all
four nevertheless record, without softening, that he swore he did not know
Jesus. Matthew's ending serves for them all: ``he began to curse and to swear
that he knew not the man. And immediately the cock crew\ldots\ And going
forth, he wept bitterly'' (Mt 26:74--75).

Of Peter at the cross the Synoptics and John say nothing. The last sight of
him in the Passion narratives is this exit in tears, and the silence that
follows it is complete: no Gospel places him at Golgotha, at the burial, or
anywhere at all until Easter morning.
